\documentclass[11pt,letterpaper]{article}
% preamble.tex — Shared preamble for all documents
\usepackage{graphicx}
\usepackage{hyperref}
\usepackage{enumitem}
\usepackage{titlesec}
\usepackage{fancyhdr}
\usepackage{setspace}
\usepackage{parskip}
\usepackage{etoolbox}
\usepackage{xspace}
\usepackage{needspace}
\usepackage{changepage}
\usepackage{geometry}
\usepackage{ragged2e}
\usepackage{environ}

% Hyperref setup
\hypersetup{
  colorlinks=false,
  pdfborder={0 0 0},
  pdftitle={The Letter to the Father},
  pdfauthor={},
}

% Load project files
% colors.tex — Color palette for The Letter to the Father
\usepackage{xcolor}

% Primary palette — muted, period-appropriate, prestige
\definecolor{kafkablack}{HTML}{1A1A1A}
\definecolor{kafkagray}{HTML}{4A4A4A}
\definecolor{kafkalight}{HTML}{8A8A8A}
\definecolor{kafkacream}{HTML}{F5F0E8}
\definecolor{kafkasepia}{HTML}{704214}
\definecolor{kafkablood}{HTML}{6B1C23}
\definecolor{kafkaink}{HTML}{2C2C3A}

% Accent colors — for package materials
\definecolor{prestigegold}{HTML}{B8860B}
\definecolor{sectionrule}{HTML}{3A3A3A}
\definecolor{pullquote}{HTML}{5A4A3A}

% metadata.tex — Project metadata for The Letter to the Father
\newcommand{\filmtitle}{The Letter to the Father}
\newcommand{\filmtitleupper}{THE LETTER TO THE FATHER}
\newcommand{\filmsubtitle}{An Original Screenplay}
\newcommand{\filmauthor}{Based on the life and writings of Franz Kafka}
\newcommand{\filmyear}{2026}
\newcommand{\filmdraft}{Revised Draft -- Public Development Example}
\newcommand{\filmdate}{October 2026 -- Example Edition}
\newcommand{\filmlogline}{In the final years of his life, a frail and tormented Franz Kafka struggles to write, love, and exist within a world of family pressure, bureaucratic humiliation, illness, and spiritual dread --- while the stories he creates begin to mirror and consume the life he can no longer control.}
\newcommand{\filmgenre}{Prestige Drama / Literary Biopic}
\newcommand{\filmrating}{R}
\newcommand{\filmlead}{Adrien Brody}
\newcommand{\filmbudgetrange}{\$12--18M}
\newcommand{\filmruntime}{Approximately 110 minutes}


% Public example status travels with every downloadable document.
\AddToHook{shipout/foreground}{\put(40,-776){\makebox(0,0)[l]{\fontsize{8}{9}\selectfont\textit{Development example. All casting, financing and awards language is illustrative; no attachments confirmed.}}}}

% typography.tex — Typography settings
\usepackage{fontspec}
\usepackage{microtype}

% Prefer repo-local metric-compatible fonts so Tectonic does not depend on
% absolute host-system font paths. Fall back to system fonts only if the vendored
% font payload is unavailable.
\IfFileExists{../../../fonts/Tinos-Regular.ttf}{
  \setmainfont{Tinos-Regular.ttf}[
    Path=../../../fonts/,
    BoldFont=Tinos-Bold.ttf,
    ItalicFont=Tinos-Italic.ttf,
    BoldItalicFont=Tinos-BoldItalic.ttf,
  ]
}{
  \IfFontExistsTF{TeX Gyre Termes}{
    \def\projectmainfont{TeX Gyre Termes}
  }{
    \IfFontExistsTF{Times New Roman}{
      \def\projectmainfont{Times New Roman}
    }{
      \def\projectmainfont{Latin Modern Roman}
    }
  }
  \setmainfont{\projectmainfont}
}

\IfFileExists{../../../fonts/Arimo-Regular.ttf}{
  \setsansfont{Arimo-Regular.ttf}[
    Path=../../../fonts/,
    BoldFont=Arimo-Bold.ttf,
    ItalicFont=Arimo-Italic.ttf,
    BoldItalicFont=Arimo-BoldItalic.ttf,
  ]
}{
  \IfFontExistsTF{TeX Gyre Heros}{
    \def\projectsansfont{TeX Gyre Heros}
  }{
    \IfFontExistsTF{Helvetica}{
      \def\projectsansfont{Helvetica}
    }{
      \def\projectsansfont{Latin Modern Sans}
    }
  }
  \setsansfont{\projectsansfont}
}

\IfFileExists{../../../fonts/Cousine-Regular.ttf}{
  \setmonofont{Cousine-Regular.ttf}[
    Path=../../../fonts/,
    BoldFont=Cousine-Bold.ttf,
    ItalicFont=Cousine-Italic.ttf,
    BoldItalicFont=Cousine-BoldItalic.ttf,
    Scale=0.95,
  ]
  \def\screenplayfont{%
    \fontspec{Cousine-Regular.ttf}[%
      Path=../../../fonts/,
      BoldFont=Cousine-Bold.ttf,
      ItalicFont=Cousine-Italic.ttf,
      BoldItalicFont=Cousine-BoldItalic.ttf,
      Scale=1.0,
    ]%
  }
}{
  \IfFontExistsTF{Courier New}{
    \def\projectmonofont{Courier New}
  }{
    \IfFontExistsTF{Courier}{
      \def\projectmonofont{Courier}
    }{
      \def\projectmonofont{Latin Modern Mono}
    }
  }
  \setmonofont{\projectmonofont}[Scale=0.95]
  \def\screenplayfont{\fontspec{\projectmonofont}[Scale=1.0]}
}

% Microtype settings
\microtypesetup{
  protrusion=true,
  expansion=false,
}

% studio_package_macros.tex — Formatting macros for the studio package

% ============================================================
% SECTION STYLING
% ============================================================
\titleformat{\section}
  {\Large\bfseries\color{kafkablack}\sffamily}
  {}
  {0em}
  {}
  [\vspace{2pt}\textcolor{sectionrule}{\rule{\textwidth}{0.5pt}}\vspace{6pt}]

\titleformat{\subsection}
  {\large\bfseries\color{kafkagray}\sffamily}
  {}
  {0em}
  {}

\titleformat{\subsubsection}
  {\normalsize\bfseries\color{kafkagray}\sffamily}
  {}
  {0em}
  {}

% ============================================================
% PULL QUOTE
% ============================================================
% Usage: \pullquote{Text here}
\NewEnviron{pullquote}{%
  \vspace{8pt}%
  \begin{adjustwidth}{0.5in}{0.5in}%
    \color{pullquote}\itshape\large%
    \BODY%
  \end{adjustwidth}%
  \vspace{8pt}%
}

% ============================================================
% HIGHLIGHT BOX
% ============================================================
\newcommand{\highlightbox}[1]{%
  \par\vspace{6pt}%
  \noindent\colorbox{kafkacream}{%
    \parbox{\dimexpr\textwidth-2\fboxsep}{%
      \vspace{4pt}#1\vspace{4pt}%
    }%
  }%
  \par\vspace{6pt}%
}

% ============================================================
% TARGET ENTRY (for director/producer lists)
% ============================================================
% Usage: \target{Name}{Why they fit}{Tier}
\newcommand{\target}[3]{%
  \needspace{4\baselineskip}%
  \subsubsection*{#1 \hfill \textnormal{\small\textit{#3}}}%
  #2\par\vspace{6pt}%
}

% ============================================================
% SCENE CARD (for showcase section)
% ============================================================
% Usage: \scenecard{Title}{Description}
\newcommand{\scenecard}[2]{%
  \needspace{5\baselineskip}%
  \vspace{4pt}%
  \noindent\textbf{\sffamily #1}\\[3pt]%
  #2\par\vspace{8pt}%
  \textcolor{sectionrule}{\rule{0.3\textwidth}{0.25pt}}%
  \vspace{8pt}%
}

% ============================================================
% TAGLINE
% ============================================================
\newcommand{\tagline}[1]{%
  \vspace{4pt}%
  \begin{center}%
    {\large\itshape\color{kafkagray} ``#1''}%
  \end{center}%
  \vspace{4pt}%
}

% ============================================================
% COMP FILM ENTRY
% ============================================================
% Usage: \compfilm{Title (Year)}{Relevance}
\newcommand{\compfilm}[2]{%
  \noindent\textbf{#1}\\%
  {\small #2}\par\vspace{6pt}%
}

% ============================================================
% RISK ENTRY
% ============================================================
% Usage: \risk{Risk}{Mitigation}
\newcommand{\risk}[2]{%
  \noindent\textbf{Risk:} #1\\%
  \textbf{Mitigation:} #2\par\vspace{8pt}%
}

% ============================================================
% PACKAGE GEOMETRY
% ============================================================
\newcommand{\packagegeometry}{%
  \newgeometry{
    left=1in,
    right=1in,
    top=1in,
    bottom=1in,
  }%
}

\geometry{left=0.75in, right=0.75in, top=0.75in, bottom=0.75in}
\pagestyle{empty}
\setlength{\parindent}{0pt}
\setlength{\parskip}{5pt}

\begin{document}

\begin{center}
{\Large\bfseries\sffamily \filmtitleupper}\\[3pt]
{\normalsize\sffamily\color{kafkagray} European Co-Production Brief}
\end{center}
\vspace{2pt}
\textcolor{sectionrule}{\rule{\textwidth}{0.5pt}}
\vspace{4pt}

\textbf{Why This Is a Natural Co-Production}\\
This is not an English-language prestige film awkwardly draped over European locations. It is materially European in subject, geography, culture, and production logic. Kafka belongs to Prague, Berlin, Vienna, and the broader Mitteleuropa imagination. The film's dramatic world is built from the actual institutions, apartments, cafés, sanatoriums, and border-crossing identities that define that region. A Czech/German/Austrian co-production is not just efficient financing. It is the correct artistic container for the movie.

\vspace{2pt}

\textbf{Territory Fit}\\
Prague is the primary production base and the emotional center of the film. Germany is essential both as production partner and as cultural market, given Kafka's language and canon status. Austria is not a financing ornament; the late-life sanatorium and final movement give it genuine narrative relevance. The project also has natural resonance across France, Italy, the UK, and prestige European arthouse territory more broadly because Kafka is already part of the cultural bloodstream there.

\vspace{2pt}

\textbf{Why the Structure Is Attractive}\\
The film is period-specific but physically contained. Repeating interior worlds, concentrated design spend, and a modest number of true set pieces make it a strong fit for a multi-territory co-production model. Czech production infrastructure, German support mechanisms, and Austrian incentives can all be used in a way that feels organic to the material rather than imposed by finance engineering.

\vspace{2pt}

\textbf{Funds and Incentive Logic}\\
Czech rebate support, German federal and regional production funding, and Austrian incentive participation all map cleanly onto the script. The combined effect can materially reduce the effective budget while preserving the prestige look on screen. Just as important, the co-production route increases cultural legitimacy with European partners, broadcasters, festivals, and sales agents who recognize that this is exactly the sort of cross-border historical prestige film the system is built to support.

\vspace{2pt}

\textbf{Why Partners Can Believe in It}\\
Adrien Brody gives the package immediate international value. Kafka gives it cultural centrality. The office storyline, father-letter engine, and Berlin movement make the film more propulsive and emotionally legible than the usual literary-biopic lane. That means co-production partners are not being asked to back a tasteful homework assignment. They are being asked to back a serious adult drama with real awards potential and clear export value.

\vspace{2pt}

\textbf{Best Partner Profile}\\
This wants co-production partners who understand prestige drama, European public-funding realities, and difficult adult material that still needs to move through festivals, sales, and awards. The ideal configuration combines Czech production muscle, German financing and market strength, Austrian final-act relevance, and a sales partner who can position the film as both culturally serious and commercially legible.

\vspace{2pt}

\textbf{What the Co-Production Route Buys Creatively}\\
Authenticity of place. Stronger access to period architecture and regional performers. Better alignment between the film's subject and its means of production. And a cleaner case to festivals and awards bodies that this is not merely a European-set prestige object, but a film made through the same cultural geography that produced its subject.

\vspace{4pt}
\textcolor{sectionrule}{\rule{\textwidth}{0.5pt}}
\vspace{2pt}
\begin{center}
{\small\color{kafkalight}\textit{A co-production not because the film needs financial patchwork, but because the material itself belongs to that structure.}}
\end{center}

\end{document}
